\section{Commercial execution platforms and open client contracts}
\label{exec23:clients-scope}

A bank's overnight causal analysis may run in an established execution
platform while its scientific meaning remains the responsibility of the
proposed library. A useful commercial comparison therefore follows the
actual boundary: submission, observation of a remote job, interpretation of
its terminal state, recovery and custody of its outputs. The open clients
supplied by SAS, Domino and Posit reveal concrete integration choices. Their
selected source revisions precede 7 October 2026. They authenticate client
and SDK behavior visible in source, not every feature of a licensed server,
a supported deployment combination, or a successful native trial.
\citep{exec23_workflow_sasctl,exec23_workflow_dominopy,
exec23_workflow_connectapi,exec23_workflow_positgo}

\subsection{Completion is one coordinate of a scientific result}
\label{exec23:clients-completion}

Let $q$ denote an immutable scientific request, $a$ a run attempt, $v$ a code
revision, $d$ a data snapshot and $e$ an execution environment. The requested
result has a declared estimand, population, information time and assumptions.
An execution record is a different object: it says which attempt ran which
inputs, whether it terminated, and which artifacts it produced. A proposed
acceptance rule can make these requirements explicit:
\begin{equation}
 A(q,a)=T(a)\,S(a)\,B(q,a,v,d,e)\,E(a)\,C(q,r_a),
 \label{exec23:clients-acceptance}
\end{equation}
where each factor is a zero--one indicator. Here $T$ means terminal, $S$
means an allowed successful execution state, $B$ means a verified binding
between request and execution identities, $E$ means the required result
artifacts $r_a$ exist, and $C$ means the result passes the stated scientific
contract. The last factor may require an explicit nonidentification or
refusal result; a numerical effect need not exist for every valid request.
These are proposed application semantics, not a claim that a vendor implements
this particular rule. A successful process exit supplies evidence for only
part of the conjunction.

This distinction matters even when all arithmetic is exact. Suppose a
funding decision has a stipulated incremental response coefficient $b_v$,
capital exposure of USD 1,000 and cost of USD 20. For two revisions that
encode different supplied assumptions, take $b_0=3/100$ and $b_1=4/100$.
Their net values are
\begin{equation}
 V_0=1000(3/100)-20=10,\qquad
 V_1=1000(4/100)-20=20.
 \label{exec23:clients-revision}
\end{equation}
Both calculations can exit successfully. Substituting the second revision
while describing the run as a replay of the first changes the answer by
USD 10. This is an authored identity example with stipulated coefficients,
not an estimated causal effect or a measured vendor failure. Recording a
code commit alone also leaves the data snapshot, configuration, dependency
versions and randomization state to be specified.
\citep{exec23_workflow_client_reference}

\subsection{Domino: versioned restart and distinct polling predicates}
\label{exec23:clients-domino}

The selected Domino Python client exposes useful controls for this boundary.
Its job submission payload includes a commit identifier, main-repository Git
reference, environment identifier and hardware selection. Its restart method
defaults to requesting the original input commit; the alternative requests
the current version. The stop method separately accepts a flag asking to
commit results. These are meaningful strengths for controlled experimentation,
but the selected client source does not establish complete data and environment
immutability or the licensed server's interpretation of every field.
\citep{exec23_workflow_dominopy}

Two polling interfaces in the same source have different visible predicates.
The legacy blocking run helper resets its consecutive status-read retry count
after a successful read. It waits for an output commit and then requires the
reported status to be \texttt{Succeeded}. The newer blocking job helper
raises on \texttt{executionStatus == Failed} and passes
\texttt{isCompleted} to the polling library as its success predicate. Its
defaults specify a five-second poll interval and a 6,000-second polling
timeout, with selected request exceptions ignored during polling.
\citep{exec23_workflow_dominopy}

These observations justify a precise question rather than an accusation of
a product defect. A supplied hypothetical response whose completion flag is
true and whose execution status is neither Failed nor Succeeded satisfies
the newer local predicate. It would fail a proposed explicit Succeeded
allowlist. This is a logical witness about the inspected expression; it does
not show that the server actually emits such a response. Resolving that
question needs the supported server state schema and its version. An adapter
should preserve both fields, the raw response, job identity and output binding
until the mapping is justified. It must also distinguish failure to observe
a job before the deadline from proof that the remote job has stopped.
\citep{exec23_workflow_client_reference,exec23_workflow_dominopy}

For the proposal, the preferred route is a narrow client wrapper that keeps
Domino's execution and revision controls and adds the scientific request
contract. Reimplementing the entire scheduler has no demonstrated advantage.
The wrapper remains a proposed integration: supported API routes, exception
mapping, actual terminal states, output custody and fault behavior still need
their own inspection and registered public or synthetic trials.

\subsection{Posit Connect: tasks, jobs and schedules have different meanings}
\label{exec23:clients-connect}

The inspected R client declares development version 0.12.1.9000. Its task
object normalizes an incoming identifier when necessary. The
\texttt{poll\_task} implementation waits for the task's
\texttt{finished} field, checks the final numeric code, raises on nonzero
code and attaches an optional result to the task object. The function exposes
a wait interval but no overall deadline parameter. This describes the local
function signature and loop, not the absence of time limits elsewhere in the
server or transport. An integrating worker should impose an explicit
observation budget and retain the remote task identity for reconciliation.
\citep{exec23_workflow_connectapi}

Job records carry additional context: content and bundle identifiers, start
and end times, execution status, exit code and, where available, cluster and
image information. The selected documentation distinguishes Active, Finished
and Finalized on newer servers and identifies missing fields on versions
earlier than 2022.10.0. The implementation also normalizes older timestamp
and identifier representations. A consumer must not treat missing historical
fields as verified empty values or infer a successful scientific result from
the label Finished. A bundle identifier is useful provenance, but its
mapping to the request's code, data and configuration still needs evidence.
\citep{exec23_workflow_connectapi}

Scheduling adds another boundary. The inspected schedule-update method is
marked experimental, converts start and next-run times and calls unversioned
schedule endpoints. Thus the R package's source commit, the server API
behavior and the schedule's information clock are separate versioning
questions. In banking, a job scheduled after a ledger cutoff must not quietly
use records first available after the intended decision time. The proposed
adapter should preserve that time explicitly and reject ambiguous timestamp
conversion. This is an application requirement, not a finding that Posit
violates it.
\citep{exec23_workflow_connectapi}

The R client is therefore a plausible integration route for R-based
estimation and reporting. Its task, job and schedule objects should remain
distinct in the library/CLI/MCP contract. Porting the client wholesale to
another language would add an API-maintenance obligation without yet proving
better scientific fidelity. A worker boundary with typed results is the
more limited candidate to test first.

\subsection{Posit Launcher: reuse the scheduler boundary with explicit state}
\label{exec23:clients-launcher}

The Launcher Go SDK addresses a different layer. Its official guide describes
a plugin connecting Workbench or Connect to a scheduler using length-prefixed
JSON over standard input and output. Bootstrap reconciles current jobs,
heartbeats detect an unresponsive plugin, and multiple Launcher instances
require a common view of job state. The guide gives Go and product-version
prerequisites. These are documented design contracts of the pinned SDK;
they are not a measured availability or recovery guarantee.
\citep{exec23_workflow_positgo}

The testing guide offers a useful structure for proposed qualification. It
separates universal invariants, product request sequences and individual
control scenarios. Examples cover nonempty job identifiers, distinct IDs for
two submissions, retrieval and filtering, terminal-state nonregression,
launch, stop, kill and cancel. Profiles allow scheduler-specific terminal
states and exit codes; unsupported suspend/resume can be skipped. An
integration must interpret the selected profile and skipped cases rather
than collapse the entire suite to one portable success claim. These tests
are documented upstream facilities; no passing execution is asserted here.
\citep{exec23_workflow_positgo}

For a bank with an accepted scheduler, adopting this bridge boundary is more
plausible than reproducing an execution platform inside a causal library.
The library would still own causal request identity and result validation.
The platform would own job dispatch and resource management. A prototype
comparison should test stale state after restart, duplicate submission,
status-stream interruption and missing output, including the cost of state
reconciliation. A failed bootstrap or missing shared state must remain
visible as an operating failure, not be translated into a zero causal effect.

\subsection{SAS: preserve experimental scope and explicit timeout outcomes}
\label{exec23:clients-sas}

The sasctl project exposes Python module and command-line access to SAS
Viya. Its selected Workflow class states that its partial implementation is
experimental and can change without notice. The comment refers to the
documentation situation in September 2022; it does not establish whether
public workflow documentation exists at the proposal cutoff. Submission
selects an enabled definition, maps supplied values to declared prompts
and includes their declared types and versions. Naive datetime values are
converted using local timezone when possible, with a UTC fallback if they
remain naive. A banking adapter should require an explicit timezone before
submission rather than let the worker host silently define the information
clock.
\citep{exec23_workflow_sasctl}

A separate score-execution method exposes another concrete control issue.
It returns the strings completed or failed on those reported states. Its
timeout test is inside the polling loop, while the loop guard itself limits
elapsed time; there is no explicit return after the loop. A supplied zero
timeout therefore skips the body and reaches Python's implicit None return.
For a positive budget, the timeout branch can be reached if time crosses
the threshold during an iteration, but a later failed loop guard can also
end the loop without that explicit timeout result. This is source-level
control-flow reasoning, not a timed experiment against SAS.
\citep{exec23_workflow_sasctl,exec23_workflow_client_reference}

The proposed wrapper should validate a positive observation budget, normalize
all completion paths into typed states, retain execution identity and logs,
and require a separately verified result binding. An unrecognized or absent
polling return is an unresolved observation, not permission to use stale
scores. Reusing the open client selectively is a candidate; depending on its
experimental workflow service without versioned qualification is not yet a
supported product commitment.

\subsection{Retries, rights and a costed prototype choice}
\label{exec23:clients-choice}

Recovery can change value when a retry repeats a charge or an externally
visible action. In an exact illustrative servicing example, suppose one
useful contact produces a stipulated USD 12 gain and every attempted
delivery incurs USD 7. If there is at most one useful gain, net value after
$n$ charged attempts is
\begin{equation}
 V(n)=12-7n,\qquad V(1)=5,\qquad V(2)=-2.
 \label{exec23:clients-retry}
\end{equation}
The sign changes solely because of the assumed duplicate cost. A repeated
analysis job need not incur this contact cost; the example concerns a
workflow that crosses into a paid or externally visible action. It motivates
separate request and attempt identities, reconciliation before resubmission
and an idempotency key at the side-effect boundary. It establishes no vendor
failure rate, bank return or economic superiority.
\citep{exec23_workflow_client_reference}

The selected connectapi license is MIT; the Launcher root license provides
an MIT-style grant unless supplied under a commercial agreement. The Domino
and sasctl client roots include Apache 2.0 texts. These source licenses are
relevant to proposed client reuse. They do not supply an entitlement to a
commercial server, a supported edition, a complete dependency-rights review
or a price. Consequently the commercial comparison should price an actual
deployment and support arrangement separately from the open wrapper.
\citep{exec23_workflow_connectapi,exec23_workflow_positgo,
exec23_workflow_dominopy,exec23_workflow_sasctl}

The practical proposal is to compare narrow wrappers around established
execution services using the same public or synthetic scientific workload.
Record installation and maintenance effort, queue and compute time, output
validation, failed attempts, recovery and any billed actions. Preserve every
requested case in the denominator. Qualify client and server versions,
explicit failure states, output custody and clock semantics before drawing
an economic conclusion. Historical Dataiku technical-version gaps and wider
market coverage remain separate survey obligations; these four inspected
profiles neither establish market completeness nor erase those limits.
